% !TEX root = ../print.tex
% Draft \S10 --- The scale-Cauchy problem: Theorem 3'
%
% Source: draft/hemigroup-causal-scale-space-kernels.md, lines 386-463.
% Chapter 10 of the blueprint; the shared counter puts every number here on the draft's own.
%
% Theorem 3'. This chapter upgrades \ref{prop:scale-evolution} from a distributional identity
% to a classical Cauchy problem on an explicit operator core, and it is the first of the three
% primed theorems.
%
% \ref{thm:scale-cauchy} carries the longest proof in the article. Its two halves use quite
% different machinery and the formalisation should treat them separately: existence is Fubini
% plus the fundamental theorem of calculus for Bochner integrals (Mathlib has both), while
% uniqueness is a SCALAR linear ODE in the Laplace variable -- deliberately not a Hille--Yosida
% argument. That choice is what keeps the chapter inside the verified core: the only [A] node
% here is \ref{prop:fixed-scale-semigroup}, which is not used by \ref{thm:scale-cauchy} at all.
% It records the semigroup reading for its own sake.
%
% Notation warning carried from ledger A11: Phillips never writes -\varphi(-A), and he asserts
% D(B) contains D(A) rather than that D(A) is a core. The two citations in A11 are genuinely
% complementary, not redundant, and the node says so.

This chapter upgrades \ref{prop:scale-evolution} from a distributional identity to a classical
Cauchy problem in the scale variable, on an explicit operator core. Throughout we work in the
canonical gauge and write
\[
  \Xz := L^1(\RR_+),
\]
identified with the causal elements of $X$; by (A3) every $\Phi_{x,y}$ restricts to $\Xz$. On
$\Xz$ let $(T_r)_{r \ge 0}$ denote the \emph{delay semigroup},
$(T_r f)(t) = f(t - r)\mathbf{1}_{t \ge r}$, a strongly continuous semigroup of positive
isometries, and define the core
\[
  \core \;:=\; \bigl\{\, f \in \Xz :\ f \text{ absolutely continuous},\ f' \in \Xz,\ f(0) = 0 \,\bigr\}.
\]

% draft: Lemma 10.1
\begin{lemma}[delay derivative]\label{lem:delay-core}
\uses{lem:convolution-representation,def:cascade-family}
$\core$ is dense in $\Xz$ and invariant under every $T_r$ and every $\Phi_{x,y}$, and for
$f \in \core$,
\[
  \lim_{h \downarrow 0} \tfrac{1}{h}\bigl(T_h f - f\bigr) = -\,f' \quad \text{in } \Xz,
  \qquad
  \bigl\|T_r f - f\bigr\|_1 \le \min\bigl(2\|f\|_1,\ r\,\|f'\|_1\bigr).
\]
\statusT
\end{lemma}

\begin{proof}
Density is standard.

\emph{Invariance under $T_r$.} $T_r f$ is absolutely continuous, vanishes on $[0, r]$, and
$(T_r f)' = T_r f'$.

\emph{Invariance under $\Phi_{x,y} = \mu_{x,y} * \cdot$.} From $f = \mathbf{1}_{[0,\infty)} * f'$
we get $\mu * f = \mathbf{1}_{[0,\infty)} * (\mu * f')$, so $\mu * f$ is absolutely continuous
with derivative $\mu * f' \in \Xz$ and value $0$ at $0$.

\emph{The limit.} On $[h, \infty)$,
$\tfrac{1}{h}(f(t-h) - f(t)) + f'(t) = -\tfrac{1}{h}\int_{t-h}^{t}\bigl(f'(\rho) - f'(t)\bigr)d\rho$,
whose $L^1$-norm tends to $0$ by continuity of translation in $L^1$; on $[0, h)$ the
contribution is $\tfrac{1}{h}\int_0^h |f| + \int_0^h |f'| \le \sup_{[0,h]}|f| + o(1) \to
|f(0)| = 0$.

\emph{The estimate.} $T_r f - f = -\int_0^r T_\rho f'\,d\rho$ as a Bochner integral --- the
integrated form of the limit --- together with $\|T_\rho f'\|_1 = \|f'\|_1$ and the trivial
bound $\|T_rf - f\|_1 \le 2\|f\|_1$.
\end{proof}

% draft: Definition 10.2
\begin{definition}[per-scale generator; Phillips form]\label{def:phillips-generator}
\uses{lem:selfdecomposable-exponents,lem:delay-core}
Let $\nu_1$ be the Lebesgue--Stieltjes measure $-dk$ on $(0,\infty)$, with $k$ taken
right-continuous and $k(\infty) = 0$, so that $\nu_1\bigl((r,\infty)\bigr) = k(r)$; and for
$x > 0$ put $\nu_x := x^{-1}$ times the pushforward of $\nu_1$ under $r \mapsto xr$, i.e.\
$\nu_x\bigl((r,\infty)\bigr) = k(r/x)/x$. For $f \in \core$ define
\[
  \varphi_x(\partial_t)\, f \;:=\; b_0\, f' \;+\; \int_0^\infty \bigl(f - T_r f\bigr)\,\nu_x(dr)
  \;=\; b_0\, f' \;+\; \frac{1}{x}\int_0^\infty \bigl(f - T_{xr} f\bigr)\,\nu_1(dr).
\]
\statusT
\end{definition}

% draft: Lemma 10.3
\begin{lemma}[properties of the generator]\label{lem:generator-properties}
\uses{def:phillips-generator,lem:delay-core,lem:memory-kernel,prop:laplace-uniqueness}
For $f \in \core$ and $x > 0$:
\begin{enumerate}
\item The integral converges absolutely in $\Xz$, with
$\|\varphi_x(\partial_t) f\|_1 \le b_0 \|f'\|_1 + \int_0^\infty \min\bigl(2\|f\|_1,\, xr\,\|f'\|_1\bigr)\,x^{-1}\nu_1(dr) < \infty$.
\item $\Lap\bigl[\varphi_x(\partial_t) f\bigr](s) = \varphi_x(s)\,\hat f(s)$ for $s > 0$, where
$\varphi_x(s) = s F'(xs)$.
\item $\varphi_x(\partial_t)$ commutes with every $\Phi_{y,z}$ on $\core$.
\item $x \mapsto \varphi_x(\partial_t) f$ is continuous from $(0,\infty)$ to $\Xz$.
\item $\kappa^{(x)} * f$ is a.e.\ equal to an absolutely continuous function with
$\bigl(\kappa^{(x)} * f\bigr)' = \varphi_x(\partial_t) f$ a.e.; that is, the Phillips form
coincides on $\core$ with the memory-kernel operator of Chapter~9.
\end{enumerate}
\statusT
\end{lemma}

\begin{proof}
(1) Combine \ref{lem:delay-core} with $\int (1 \wedge r)\,\nu_x(dr) < \infty$, which follows
from $\int_0^1 k < \infty$ and $k(1) < \infty$ by integration by parts.

(2) Termwise: $\Lap[f'] = s\hat f$, using $f(0) = 0$, and $\Lap[f - T_rf] = (1 - e^{-sr})\hat f$.
Fubini is justified by (1), and $b_0 s + \int(1 - e^{-sr})\nu_x(dr) = \varphi_x(s)$ by
\ref{lem:memory-kernel} --- the two Bernstein representations of $\varphi_x$ agree.

(3) Convolutions commute; pull $\Phi_{y,z}$ through the Bochner integral.

(4) In the dilated form of \ref{def:phillips-generator} the integrand
$\tfrac1x(f - T_{xr}f)$ is, for $x$ in a compact subset of $(0,\infty)$, continuous in $x$ for
each $r$ by strong continuity of $T$, and dominated by $C\min(\|f\|_1,\, r\|f'\|_1)$, which is
$\nu_1$-integrable. Apply dominated convergence; the drift term is trivial.

(5) Both $\varphi_x(\partial_t)f \in L^1$ and $\kappa^{(x)} * f \in L^1_{\mathrm{loc}}$ are
causal with finite Laplace transforms. By (2) and \ref{lem:memory-kernel} the causal function
$V(t) := \int_0^t \bigl(\varphi_x(\partial_t)f\bigr)(\rho)\,d\rho - (\kappa^{(x)} * f)(t)$ has
transform $\varphi_x(s)\hat f(s)/s - F'(xs)\hat f(s) = 0$ for all $s > 0$, so $V = 0$ a.e.\ by
\ref{prop:laplace-uniqueness}. That is the claim; the value $(\kappa^{(x)}*f)(\zp) = 0$
follows from $\kappa^{(x)}([0,t]) \to b_0$ and $\sup_{[0,t]}|f| \to 0$.
\end{proof}

% draft: Theorem 10.4
\begin{theorem}[Theorem 3$'$: the scale-Cauchy problem]\label{thm:scale-cauchy}
\uses{thm:main-characterization,lem:delay-core,def:phillips-generator,lem:generator-properties,prop:scale-evolution,def:cascade-family,prop:laplace-uniqueness}
Let $(\Phi_{x,y})$ be as in \ref{thm:main-characterization}, in the canonical gauge, and let
$f \in \core$. Set $u(\cdot, x) := \Phi_{0,x} f$. Then:
\begin{enumerate}
\item $u(\cdot, x) \in \core$ for every $x \ge 0$;
\item $x \mapsto u(\cdot, x)$ belongs to
$C\bigl([0,\infty); \Xz\bigr) \cap C^1\bigl((0,\infty); \Xz\bigr)$ and satisfies
\[
  \partial_x u(\cdot, x) \;=\; -\,\varphi_x(\partial_t)\,u(\cdot, x) \quad (x > 0),
  \qquad u(\cdot, 0) = f,
\]
with $\partial_x$ a norm-derivative in $\Xz$;
\item $u$ is the \emph{unique} solution in the class of
$v \in C([0,\infty); \Xz) \cap C^1((0,\infty); \Xz)$ with $v(\cdot, x) \in \core$ for $x > 0$,
satisfying the equation for $x > 0$ and $v(\cdot, 0) = f$.
\end{enumerate}
\statusT\quad\emph{(the hemigroup analogue of \sscite{fagerstrom2005temporal}, Thm.~3, which
is recovered as the homogeneous case in \ref{rem:recovering-thm3}. Proved here, not cited.)}
\end{theorem}

\begin{proof}
(1) is the $\Phi$-invariance of $\core$ (\ref{lem:delay-core}), and continuity at every
$x \ge 0$ is (A7).

(2) Write $v(\cdot, x) := \varphi_x(\partial_t)\, u(\cdot, x)$. By
\ref{lem:generator-properties}(3), $v(\cdot, x) = \Phi_{0,x}\,\varphi_x(\partial_t) f$, so by
\ref{lem:generator-properties}(4), contractivity of $\Phi_{0,x}$, and (A7),
\[
  \|v(\cdot,x) - v(\cdot,x')\|_1
  \;\le\; \bigl\|\varphi_x(\partial_t)f - \varphi_{x'}(\partial_t)f\bigr\|_1
  + \bigl\|(\Phi_{0,x} - \Phi_{0,x'})\,\varphi_{x'}(\partial_t)f\bigr\|_1
  \;\longrightarrow\; 0,
\]
so $x \mapsto v(\cdot, x)$ is continuous from $(0,\infty)$ to $\Xz$.

By \ref{lem:generator-properties}(5) applied to $u(\cdot,x) \in \core$, each
$\kappa^{(x)} * u(\cdot, x)$ is absolutely continuous with derivative $v(\cdot, x)$, vanishing
at $0$; hence for $0 < a \le b$, Fubini gives
\[
  \int_a^b \bigl(\kappa^{(x)} * u(\cdot, x)\bigr)(t)\,dx
  \;=\; \int_0^t \int_a^b v(\rho, x)\,dx\,d\rho ,
\]
so the distributional $\partial_t$ in \ref{prop:scale-evolution}(2) evaluates to the Bochner
integral $\int_a^b v(\cdot, x)\,dx$:
\[
  u(\cdot, b) - u(\cdot, a) \;=\; -\int_a^b v(\cdot, x)\,dx .
\]
Since the integrand is continuous in $x$, the fundamental theorem of calculus for Bochner
integrals yields $u \in C^1((0,\infty); \Xz)$ with $\partial_x u = -v$, which is the equation.

(3) Let $v$ be any solution in the stated class and fix $s > 0$. Then $\eta(x) := \hat v(s, x)$
is continuous on $[0,\infty)$ and $C^1$ on $(0,\infty)$, with
\[
  \eta'(x) = \Lap[\partial_x v](s) = -\Lap[\varphi_x(\partial_t) v](s) = -\varphi_x(s)\,\eta(x)
\]
by \ref{lem:generator-properties}(2); differentiation and $\Lap$ exchange because
$|\hat g(s)| \le \|g\|_1$. Solving this scalar linear ODE on $[\varepsilon, x]$ and letting
$\varepsilon \downarrow 0$, using continuity at $0$, gives $\eta(x) = \hat f(s)\, e^{-F(xs)}$.
Thus $\hat v(s, x) = \hat u(s, x)$ for all $s > 0$, and $v = u$ by
\ref{prop:laplace-uniqueness}.
\end{proof}

% draft: Proposition 10.5
\begin{proposition}[fixed-scale semigroups; core property]\label{prop:fixed-scale-semigroup}
\uses{lem:potential-kernel,lem:delay-core,def:phillips-generator}
For fixed $x > 0$, the operator $-\varphi_x(\partial_t)$ with domain $\core$ is closable, and
its closure generates the strongly continuous contraction semigroup on $\Xz$ of convolutions
by the probability measures with Laplace transforms $e^{-\tau \varphi_x(s)}$, $\tau \ge 0$ ---
the semigroup subordinate to the delay semigroup via $\varphi_x$. Moreover $\core$ is a core
for that generator.
\statusA\quad\ledger{A11}\quad\emph{(Phillips' subordination theorem,
\sscite{phillips1952generation}, Thm.~4.3 with the hypotheses in its equation~(40); and the
core criterion, \sscite{engel2000one}, Prop.~II.1.7. Cited, not reproven here.
\textbf{Assignment.} \ledger{A11} carries two complementary halves and neither suffices alone.
Phillips gives the form of the subordinate generator; he asserts that its domain \emph{contains}
$\core$, not that $\core$ is a core for it, which is precisely the gap Engel--Nagel closes.
Two things the citation does not carry, both \textbf{[T]}: that $\core$ is dense and
semigroup-invariant, proved in \ref{lem:delay-core}, and that $\varphi_x$ is a Bernstein
function without killing term, proved in \ref{lem:potential-kernel}. A notational caution
recorded in the ledger: Phillips writes neither $-\varphi(-A)$ nor anything in that notation;
his displayed form becomes it only after setting the drift and killing parameters to zero.
Note finally that \ref{thm:scale-cauchy} does \emph{not} use this proposition --- the Cauchy
problem is solved by a scalar ODE, not by generation theory, so this node adds a reading
rather than a dependency.)}
\end{proposition}

\begin{proof}
$\varphi_x$ is a Bernstein function without killing term by \ref{lem:potential-kernel}, so
Phillips' subordination theorem applies to the delay semigroup, whose generator extends
$f \mapsto -f'$ on $\core$ (\ref{lem:delay-core}): the subordinate generator is the closure of
the Phillips form on the generator's domain. That $\core$ is a core follows from the cited
criterion, since $\core$ is dense and invariant under the subordinate semigroup, which
consists of convolutions (\ref{lem:delay-core}).
\end{proof}

% draft: Example 10.6
\begin{example}[the Gamma family: time-recursive form]\label{ex:gamma-recursive}
\uses{prop:gamma-family,def:phillips-generator}
For $F(s) = \gamma\log(1+s)$ we have $\nu_x(d\rho) = \gamma x^{-2} e^{-\rho/x}\,d\rho$, an
exponential memory, and $\varphi_x(s) = \gamma s/(1 + xs)$, so
\[
  \varphi_x(\partial_t)\,u \;=\; \gamma\,\partial_t\bigl(1 + x\,\partial_t\bigr)^{-1} u
  \;=\; \gamma\, w', \qquad \text{where } w \text{ solves } w + x\,w' = u,\ \ w(0) = 0 .
\]
The auxiliary variable $w$ is computed by one causal first-order ODE, so the scale step is
realized by a single recursive low-pass filter --- the strongest possible form of the
time-recursivity requirement of \sscite{fagerstrom2005temporal}, \S5: the observer embodies
its past in one first-order state per scale.
\end{example}

% draft: Remark 10.7
\begin{remark}[recovering Theorem 3 of 2005]\label{rem:recovering-thm3}
\uses{thm:scale-cauchy,prop:stable-family,rem:gauge-freedom}
For $F(s) = s^\alpha$ with $0 < \alpha < 1$ we have
$\nu_x\bigl((\rho,\infty)\bigr) = \alpha x^{\alpha-1}\rho^{-\alpha}/\Gamma(1-\alpha)$ and
$b_0 = 0$, so $\varphi_x(\partial_t) = \alpha x^{\alpha - 1} D^{\alpha}_{t,+}$ with
$D^{\alpha}_{t,+}$ the Riemann--Liouville/Marchaud derivative, and \ref{thm:scale-cauchy} reads
$\partial_x u = -\alpha x^{\alpha-1} D^{\alpha}_{t,+} u$. Under the reparametrization
$\tau := x^{\alpha}$ this is exactly
\[
  \partial_\tau u \;=\; -\,D^{\alpha}_{t,+}\, u,
\]
Theorem~3 of \sscite{fagerstrom2005temporal} verbatim --- identifying the scale parameter
$\tau$ of that paper as the $\alpha$-gauge of \ref{rem:gauge-freedom} applied to the canonical
scale coordinate.
\end{remark}
